\chapter{Conclusions}
\label{chap:conclusion}

\section{Contributions and main findings}

Accounting for uncertainty in the estimated standard error can improve
tests of equal mutual information (MI), but does not by itself ensure
accurate false-positive rates. Across the populations studied, the proposed
Expanded Welch--Satterthwaite test (Expanded Welch) brings some excessive
false-positive rates closer to the 5\% target, while lowering rates already
below that target. Its benefit therefore depends on the errors present in
the estimated MI difference and standard error, as well as the rejection
cutoff.

The thesis derives the adjustment by calculating how changes in a
probability table affect its estimated pointwise-MI variance, including
changes in the cell probabilities and their row and column totals. The
result gives a first-order approximation to the sampling variability of
that variance estimate, which is used to estimate degrees of freedom for
a Student \(t\) reference distribution. The Expanded Welch test retains
the MI difference and standard error used by the Wald test, so the two
methods share any errors in those estimates. The simulations examine both
the approximations underlying the adjustment and its consequences for
false-positive rates and power.

The online-shopping application illustrates the comparison on observed
data. Both tests reject equal month--purchase MI between new and returning
visitors, with a higher corrected MI estimate among returning visitors.
Neither test provides evidence of a difference in weekend--purchase MI.

\subsection{RQ1: Accuracy of variance approximations}

The first-order approximations become more accurate as samples grow in
the studied populations, but larger tables require more observations.
At smaller sample sizes, the population pointwise-MI variance can
substantially understate the average of its estimates across repeated
samples. This error can lower the predicted degrees of freedom even when
the variation of those estimates around their average is predicted fairly
accurately. Substituting sample proportions for population probabilities
introduces further error and can reverse the direction of the sampling-variance
prediction error. Accurate degrees of freedom therefore require adequate
predictions of both the average variance estimate and its variation between
samples.

Near independence between the two variables, sampling noise produces
apparent association that the first-order approximation does not fully
capture. At fixed sample size, the population first-order degrees of
freedom approach zero as MI decreases, while the average sample-based
degrees of freedom remain positive in the experiment. These are different
predictions because the sample-based calculation retains the apparent
association in the observed tables. The additional sampling variation
examined in Appendix~\ref{sec:independence-boundary} explains this
separation and shows why the population first-order degrees of freedom
alone give an inadequate account of behaviour near independence.

\subsection{RQ2: False-positive rates under equal MI}

False-positive behaviour changes with shared MI, table size, category
frequencies, and sample size. Across the null grid, both tests exhibit
ranges of excessive rejection and ranges of conservative rejection, where
false-positive rates fall below 5\%. The Expanded Welch test lowers
rejection in both types of range, so a lower rate does not consistently
mean a more accurate test. The diagnostic comparisons identify different
causes of these departures: an underestimated standard error can inflate
rejection, an oversized standard error can suppress it, and unequal biases
in the two MI estimates can shift their difference away from zero.
At larger sample sizes, the methods become close, although both can remain
conservative in the weak-association, larger-table settings.

Selected settings illustrate the extent of both improvement and remaining
failure. For uniform \(4\times4\) tables at shared MI 0.02 nats and
20 observations per population, the false-positive rate falls from 8.35\%
with the Wald test to 4.90\% with the Expanded Welch test. For
different-skew \(16\times16\) tables at shared MI 0.005 nats and
1,000 observations per population, it falls only from 36.38\% to 34.38\%.
The latter failure persists when every test calculation succeeds. Together
with the conservative ranges, these results show that the larger cutoff
can correct some excessive rejection but cannot generally compensate for
inaccurate estimates of the MI difference or standard error.

\subsection{RQ3: Power across MI differences}

At the same 5\% significance level, the Expanded Welch test cannot reject
where the Wald test would not, because it uses the same statistic with a
larger cutoff. Its power therefore cannot exceed the Wald test's power
under these testing rules. Unsuccessful Expanded Welch calculations also
reduce reported power at some of the smallest sample sizes. The power
comparison must be read alongside the null results: reducing rejection is
beneficial where false positives are excessive, but can also make genuine
differences harder to detect. The experiments do not compare power after
matching the tests' actual false-positive rates.

The methods become closer with larger samples: at 10,000 observations per
population, their power difference is below 0.32 percentage points across
all tested settings. A separate difficulty in sparse settings is
that bias can reverse which population appears to have greater MI, even
when a test rejects equality. Bias can also oppose the true difference,
so increasing the population MI gap initially moves the average estimated
difference towards zero and reduces rejection. Detecting unequal MI and
estimating the direction of that difference must therefore be considered
together.

\section{Limitations}

The derivation assumes independent categorical observations within each
population, independent samples between populations, a fixed set of
categories, and positive probability in every cell. The degrees-of-freedom
calculation also requires positive pointwise-MI variances and positive
variances of the observation-level sensitivities. Exact independence is
treated separately because the first-order sampling term vanishes there.
Paired observations, continuous variables, and a growing number of
categories require extensions to the present method.

The experiments use one population construction, which increases the
probability of matching categories while preserving category frequencies,
with matching row and column frequency profiles within each table. The tested table sizes, MI values, and
margin regimes therefore cover a restricted range of association patterns.
The power study fixes one population's MI and increases the other's, while
all main simulations use equally sized samples. With equal sample sizes,
the leading bias corrections cancel from the MI difference, leaving their
performance with unequal sample sizes untested. Different residual biases
can still remain in the two MI estimates.

The main empirical benchmark is the Wald test, with ordinary Welch
degrees of freedom examined at three illustrative null settings. The
results therefore do not establish the relative performance of Expanded
Welch against resampling procedures or tests based on other MI estimators.
Likewise, the sample sizes at which the methods become close in these
experiments do not establish a threshold that guarantees accurate
inference across populations.

\section{Future work}

The remaining failures make improving the estimated MI difference and
standard error a priority. Evaluating changes to these estimates
separately would help determine which errors are reduced and how the
changes affect false-positive rates, power, and the estimated direction of
the MI difference. Comparisons with alternative MI estimators and
resampling procedures would also show whether those methods address the
failures observed here.

Retaining terms omitted by the first-order approximation could improve
predictions of the average variance estimate and its sampling variability,
and hence the degrees-of-freedom calculation. The near-independence
analysis provides a starting point for retaining variation from sampling
noise when association is weak. Any resulting test must still allow
populations with different probability tables to satisfy the null
hypothesis when their MI values are equal.

Broader experiments should examine other association patterns, different
row and column frequency profiles, unequal sample sizes, and other starting MI
values and directions of change. These comparisons would help establish
how widely the present findings apply, particularly when rare category
combinations remain poorly observed.

\section{Overall conclusion}

The thesis provides a derived adjustment for uncertainty in the estimated
standard error and explains why its benefit varies across the tested
populations. The adjustment reduces some excessive false-positive rates,
but does not provide a general improvement over the Wald test in small or
sparse samples. Reliable MI comparisons also require a difference estimate
with small bias relative to its sampling spread and a standard error that
describes that spread. Improving these estimates is the main priority
emerging from the remaining failures.
